% os-reconstruction.tex — Chapter on Osterwalder–Schrader
% reconstruction of (E_+, <,>_OS, H_infty, Omega_0, H_op) from mu_infinity.
%
% Linked code / Lean:
%   - lean/YMGap/Continuum/OSReconstruction.lean (Lean type signature
%     exposing the OS-reconstruction quintuple)
%   - \Cref{ch:os-reconstruction} of work_plan.md.

\chapter{Osterwalder--Schrader Reconstruction
  for \texorpdfstring{$\muinf$}{mu\_infinity}:
  Hilbert Space and Hamiltonian}
\label{ch:os-reconstruction}

\noindent\emph{From the OS-axiomatized $\muinf$ to the OS quintuple
$(\Ecalp, \IprodOS{\cdot}{\cdot}, \Hphys, \Omegavac, \Hop)$ ready for
spectral analysis.}

We apply the Osterwalder--Schrader reconstruction theorem to the
continuum $\SUtwo$ Wilson measure $\muinf$ of \Cref{ch:continuum-measure} (verified
to satisfy the OS axioms in \Cref{ch:os-axioms},
Chapter~\ref{ch:os-axioms}) to construct the physical Hilbert space
$\Hphys$ and the self-adjoint Hamiltonian $\Hop \ge 0$ ready for the
spectral analysis of the mass-gap chapter (Chapter~\ref{ch:mass-gap-theorem}).

The construction has four canonical steps: (i)~assemble the algebra
$\Ecalp$ of cylindrical test functions supported on the positive-
time half-space; (ii)~define the Schwinger inner product
$\IprodOS{F}{G} := \int (\theta F)^* G \, d\muinf$ on $\Ecalp$, which
is positive semidefinite by OS2; (iii)~complete the quotient
$\Ecalp / \Ncal$ by its null space $\Ncal$ to obtain $\Hphys$;
(iv)~the time-translation $\Tt$ on $\Ecalp$ extends to a strongly
continuous contraction semigroup on $\Hphys$, and Stone's theorem
gives the self-adjoint generator
$\Hop \ge 0$ such that $\Tt = e^{-t\Hop}$.

Following the work-plan's explicit open-question kickoff for 8.2,
we deliver the construction at the \textbf{paper level} (full
analytical content here), with a parallel \textbf{Lean-formal-
statement-only} type signature module
\codepath{lean/YMGap/Continuum/OSReconstruction.lean} exposing the
data of $(\Hphys, \Hop, \Omegavac)$ that downstream the mass-gap chapter
references symbolically.  The full Lean formalization (Stone's
theorem, spectral theorem for self-adjoint operators) is a major
undertaking deferred to a parallel constructive-analysis layer
(a constructive-reals layer+).

\section{Standing setup}
\label{sec:os-reconstruction-setup}

We import the OS axiom verification of \Cref{ch:os-axioms}: the measure $\muinf$ on $\Ominf$ satisfies
OS0~(regularity), OS1~(hypercubic translation and $H_4$ invariance,
contingent $O(4)$ via the restoration argument),
OS2~(reflection positivity), OS3~(internal gauge and discrete
symmetry), and OS4~(cluster property with explicit exponential rate
$\sighinf \ge \varepsilon > 0$ via the $\sighinf$-positivity theorem).

The continuum configuration space $\Ominf$ is compact Hausdorff
Polish (\Cref{ch:continuum-measure} Definition~2.4); the
cylindrical test-function algebra $\Fcal_\infty = \bigcup_k \Fcal_k$
is dense in $C(\Ominf)$ by Stone--Weierstrass.

\subsection{Notation: time direction and half-space algebra}
\label{subsec:half-space}

Fix the 4-direction $x^0$ as Euclidean time (cf.\ \texttt{os-axioms.tex}
\S1.2).  The \emph{positive-time half-space} is
$H^+ := \{ x \in \Rfour : x^0 \ge 0 \}$, with corresponding link set
$E_k^+ = \{ \ell \in E_k : \ell \subseteq H^+ \}$ at finite scale
$k$ and limit $E_\infty^+$ at the continuum.

The \emph{positive-half-space test function algebra} is
\begin{equation*}
  \Ecalp \;:=\;
  \bigl\{ F \in \Fcal_\infty : F \text{ depends only on link variables in } E_\infty^+ \bigr\}.
\end{equation*}
By the projective-limit construction, $F \in \Ecalp$ iff $F$ factors
as $F = f \circ \pi_k$ with $f \in C(\Omk)$ depending only on
$E_k^+$ link variables, for some $k$.

The time reflection $\theta \colon \Ominf \to \Ominf$ of
\texttt{os-axioms.tex} \S1.2 maps $E_\infty^+$ onto $E_\infty^-$ and
reverses the orientation of time-pointing links; pulled back to
test functions, $\theta F \in \Ecal_-$ (the analogous negative-time
algebra).

\subsection{Time-translation action}
\label{subsec:time-translation}

For each $t \ge 0$, the lattice time-translation $\tau_t$ acts on
$E_\infty^+$ by $x^0 \mapsto x^0 + t$; this restricts to a map
$\tau_t \colon E_\infty^+ \to E_\infty^+$ (positive-time
half-space is preserved by forward time-translation).

Define the pullback $\Tt \colon \Ecalp \to \Ecalp$ by
$(\Tt F)(\omega) := F(\tau_t^{-1} \omega)$.  By OS1
(\texttt{os-axioms.tex} Theorem~4.2),
$\int (\Tt F) \, d\muinf = \int F \, d\muinf$.

\section{The Schwinger inner product}
\label{sec:schwinger}

\begin{definition}[Schwinger inner product]
\label{def:schwinger}
The \emph{Schwinger inner product} on $\Ecalp$ is
\begin{equation*}
  \IprodOS{F}{G} \;:=\; \int_{\Ominf} (\theta F)^* \, G \, d\muinf
  \qquad (F, G \in \Ecalp).
\end{equation*}
\end{definition}

\begin{proposition}[Schwinger form is positive semidefinite Hermitian]
\label{prop:schwinger-psd}
For every $F, G \in \Ecalp$:
\begin{enumerate}
  \item $\IprodOS{F}{G} = \overline{\IprodOS{G}{F}}$  (Hermitian symmetry).
  \item $\IprodOS{F}{F} \ge 0$  (positive semidefiniteness).
  \item $\IprodOS{\cdot}{\cdot}$ is sesquilinear: conjugate-linear in
    the first argument, linear in the second.
\end{enumerate}
\end{proposition}

\begin{proof}
\emph{(i) Hermitian symmetry}: $\overline{\IprodOS{G}{F}} =
\overline{\int (\theta G)^* F \, d\muinf} = \int F^* (\theta G) \, d\muinf$.
Apply the reflection $\theta$ to the integration variable (which
preserves $\muinf$ by OS3 reflection symmetry, or equivalently by
the definition of $\theta$ as a measurable involution with
$\theta_* \muinf = \muinf$): $= \int (\theta F)^* G \, d\muinf =
\IprodOS{F}{G}$.

\emph{(ii) Positive semidefiniteness}: this is exactly OS2 of
\texttt{os-axioms.tex} Theorem~5.2.

\emph{(iii) Sesquilinearity}: linear in $G$ and conjugate-linear in
$F$ by definition.
\end{proof}

\begin{remark}[Schwinger form is not strictly positive definite]
\label{rem:schwinger-not-pd}
$\IprodOS{F}{F} = 0$ can hold for $F \ne 0$: any $F$ for which
$\theta F$ is uncorrelated with $F$ under $\muinf$ has zero
self-inner-product.  The quotient by the null space (next section)
is what converts the semidefinite form into a definite inner
product on the physical Hilbert space.
\end{remark}

\section{The physical Hilbert space $\Hphys$}
\label{sec:hilbert-space}

\begin{definition}[Null subspace]
\label{def:null-subspace}
The \emph{Schwinger null subspace} is
\begin{equation*}
  \Ncal \;:=\; \{ F \in \Ecalp : \IprodOS{F}{F} = 0 \}.
\end{equation*}
By the Cauchy--Schwarz inequality on the positive-semidefinite
Hermitian form (Proposition~\ref{prop:schwinger-psd}), $\Ncal$ is a
\emph{linear subspace} of $\Ecalp$ and $F \in \Ncal$ iff
$\IprodOS{F}{G} = 0$ for every $G \in \Ecalp$.
\end{definition}

\begin{definition}[Pre-Hilbert quotient]
\label{def:pre-hilbert-quotient}
The quotient $\Ecalp / \Ncal$ inherits a non-degenerate
positive-definite Hermitian inner product
\begin{equation*}
  \Iprod{[F]}{[G]} \;:=\; \IprodOS{F}{G}
\end{equation*}
that is well-defined (independent of representatives) by the
null-subspace characterization.
\end{definition}

\begin{theorem}[The physical Hilbert space]
\label{thm:physical-hilbert-space}
The \emph{physical Hilbert space} of the OS reconstruction is the
metric completion
\begin{equation*}
  \Hphys \;:=\; \overline{\,\Ecalp / \Ncal\,}^{\,\|\cdot\|_{\mathrm{OS}}}
\end{equation*}
with the inner product extended by continuity.  $\Hphys$ is a complex
Hilbert space.
\end{theorem}

\begin{proof}
Apply the standard completion of a non-degenerate inner-product
space to a Hilbert space.  Separability of $\Hphys$ follows from
separability of $\Ecalp$ in the sup norm
(consequence of the Polish-space structure of $\Ominf$).
\end{proof}

\begin{definition}[Vacuum vector]
\label{def:vacuum}
The \emph{vacuum vector} $\Omegavac \in \Hphys$ is the equivalence
class of the constant function $1 \in \Ecalp$:
\begin{equation*}
  \Omegavac \;:=\; [\,1\,] \in \Hphys.
\end{equation*}
By OS0 + $\muinf$ a probability measure,
$\IprodOS{1}{1} = \int 1 \, d\muinf = 1$, so $\|\Omegavac\| = 1$ and
$\Omegavac \ne 0$.
\end{definition}

\section{The contraction semigroup}
\label{sec:semigroup}

The time-translation $\Tt$ on $\Ecalp$ descends to the quotient
$\Ecalp / \Ncal$ and extends to $\Hphys$ as a one-parameter
strongly continuous contraction semigroup.

\begin{proposition}[$\Tt$ preserves the null space]
\label{prop:T-preserves-null}
For every $t \ge 0$ and every $F \in \Ncal$, $\Tt F \in \Ncal$.
\end{proposition}

\begin{proof}
$\IprodOS{\Tt F}{\Tt F} = \int (\theta \Tt F)^* \, \Tt F \, d\muinf$.
The reflection $\theta$ commutes with $\Tt$ up to a sign on
time-pointing links (which is absorbed into $\theta$'s definition):
$\theta \Tt = \Tt \theta$ on $\Ecalp$.  Hence
$(\theta \Tt F)^* = (\Tt \theta F)^* = \Tt (\theta F)^*$.

The time-translation $\Tt$ preserves $\muinf$ by OS1
(\texttt{os-axioms.tex} Theorem~4.2), so
$\int (\Tt (\theta F)^*) \, (\Tt F) \, d\muinf =
 \int (\theta F)^* F \, d\muinf =
 \IprodOS{F}{F} = 0$.

Therefore $\IprodOS{\Tt F}{\Tt F} = 0$, i.e.\ $\Tt F \in \Ncal$.
\end{proof}

By Proposition~\ref{prop:T-preserves-null}, $\Tt$ descends to a
linear map $\overline{\Tt} \colon \Ecalp / \Ncal \to \Ecalp / \Ncal$.

\begin{proposition}[$\overline{\Tt}$ is a contraction]
\label{prop:T-contraction}
For every $t \ge 0$ and every $[F] \in \Ecalp / \Ncal$,
$\|\overline{\Tt} [F]\| \le \|[F]\|$.
\end{proposition}

\begin{proof}[Proof sketch]
This is the classical Glimm--Jaffe contraction estimate.  Write
$F = F^+ + F^-$ where $F^+$ depends only on links with
$x^0 \ge t/2$ and $F^-$ depends only on links with
$x^0 \le t/2$ (refinement of $\Ecalp$ via the spatial
half-space at midheight $t/2$).  Then $\Tt F^+ \in \Ecal^{++}$ (the
quadruple half-space, $x^0 \ge 3t/2$) and the inner product
$\IprodOS{\Tt F}{\Tt F}$ factors through the spatial reflection at
$t/2$, which by a second application of reflection positivity
contracts the norm:
$\|\overline{\Tt} [F]\|^2 \le \|[F^+]\|^2 + \|[F^-]\|^2$.

The full argument is given in Glimm--Jaffe \S6.1 (the
``positivity-implies-contraction'' theorem); see Osterwalder--Schrader
1973 for the original.
\end{proof}

\begin{theorem}[Strongly continuous contraction semigroup]
\label{thm:semigroup}
$\overline{\Tt}$ extends by continuity to a strongly continuous
one-parameter contraction semigroup $\Ut \colon \Hphys \to \Hphys$
($t \ge 0$) satisfying:
\begin{enumerate}
  \item $U(0) = \mathrm{Id}_{\Hphys}$.
  \item $U(t + s) = U(t) \, U(s)$ for all $s, t \ge 0$.
  \item $\|\Ut\|_{\mathrm{op}} \le 1$ for every $t \ge 0$.
  \item $t \mapsto \Ut$ is strongly continuous: for every
    $\psi \in \Hphys$, $\lim_{t \to 0^+} \|\Ut \psi - \psi\| = 0$.
  \item $\Ut \Omegavac = \Omegavac$ for every $t \ge 0$  (vacuum
    is time-translation invariant).
\end{enumerate}
\end{theorem}

\begin{proof}
Items 1--3 are inherited from $\Tt$ on $\Ecalp$
(Proposition~\ref{prop:T-contraction}).  Item~5 follows from
$\Tt \cdot 1 = 1$ at the level of $\Ecalp$.  Item~4 (strong
continuity) follows from the dense-subspace approximation:
$\Ut \to \mathrm{Id}$ pointwise on the dense $\Ecalp / \Ncal$ by
continuity of $\muinf$-expectations in their integrand
(consequence of dominated convergence on the compact $\Ominf$).
\end{proof}

\section{Stone's theorem and the Hamiltonian $\Hop$}
\label{sec:stone}

\begin{theorem}[Stone's theorem for self-adjoint semigroups]
\label{thm:stone-semigroup}
Let $\Ut$ ($t \ge 0$) be a strongly continuous one-parameter
contraction semigroup on a complex Hilbert space $\Hphys$.  Then
there exists a unique self-adjoint nonnegative operator
$\Hop \ge 0$ on $\Hphys$ such that
\begin{equation*}
  \Ut \;=\; e^{-t \Hop} \qquad \text{for every } t \ge 0.
\end{equation*}
The domain $\mathcal{D}(\Hop) = \{ \psi : t \mapsto \Ut \psi \text{ is}\ C^1 \text{ at } 0 \}$
and $\Hop \psi = -\frac{d}{dt}\Big|_{t=0^+} \Ut \psi$ for $\psi \in \mathcal{D}(\Hop)$.
\end{theorem}

This is the classical Hille--Yosida + Stone construction; see
Reed--Simon \emph{Methods of Modern Mathematical Physics} Vol.~II,
Theorems X.47a and X.49.  For self-adjoint contraction semigroups
(the OS case), self-adjointness of $\Hop$ is automatic from
$\Ut^* = \Ut$ (a consequence of OS2 reflection symmetry of the
Schwinger form; cf.\ Proposition~\ref{prop:schwinger-psd}(i)
applied to $\Ut$ on $\Ecalp / \Ncal$).

\begin{theorem}[The OS Hamiltonian $\Hop$]
\label{thm:OS-hamiltonian}
There exists a unique self-adjoint nonnegative operator
$\Hop \ge 0$ on $\Hphys$ such that the time-translation semigroup
of Theorem~\ref{thm:semigroup} is $\Ut = e^{-t \Hop}$.  The vacuum
$\Omegavac$ is in the kernel: $\Hop \Omegavac = 0$.
\end{theorem}

\begin{proof}
Apply Stone's theorem (Theorem~\ref{thm:stone-semigroup}) to $\Ut$
of Theorem~\ref{thm:semigroup}; self-adjointness uses the
contraction-semigroup-is-symmetric refinement of OS reflection
positivity (Reed--Simon Vol.~II Thm.~X.47a).  Vacuum-annihilation:
$\Ut \Omegavac = \Omegavac$ for every $t$ implies
$\Hop \Omegavac = -\frac{d}{dt}\big|_{t=0^+} \Omegavac = 0$.
\end{proof}

\section{The OS reconstruction package}
\label{sec:package}

\begin{definition}[OS reconstruction quintuple]
\label{def:os-reconstruction}
The \emph{Osterwalder--Schrader reconstruction} of $(\Ominf, \muinf)$
is the quintuple
\begin{equation*}
  \mathrm{Rec}(\muinf) \;:=\; \bigl( \Hphys, \; \Omegavac, \; \Ut, \; \Hop, \; \mathcal{D}(\Hop) \bigr)
\end{equation*}
consisting of:
\begin{enumerate}
  \item The physical Hilbert space $\Hphys$ (Theorem~\ref{thm:physical-hilbert-space}).
  \item The vacuum vector $\Omegavac \in \Hphys$ with
    $\|\Omegavac\| = 1$ (Definition~\ref{def:vacuum}).
  \item The strongly continuous contraction semigroup $\Ut$
    (Theorem~\ref{thm:semigroup}).
  \item The self-adjoint Hamiltonian $\Hop \ge 0$ with
    $\Hop \Omegavac = 0$ (Theorem~\ref{thm:OS-hamiltonian}).
  \item The dense domain $\mathcal{D}(\Hop) \subseteq \Hphys$.
\end{enumerate}
\end{definition}

\begin{theorem}[Existence and uniqueness of $\mathrm{Rec}(\muinf)$]
\label{thm:reconstruction-existence}
For the continuum SU(2) Wilson measure $\muinf$ satisfying the OS
axioms (\texttt{os-axioms.tex} \S2), the reconstruction
$\mathrm{Rec}(\muinf)$ exists and is unique up to unitary equivalence
of Hilbert spaces.
\end{theorem}

\begin{proof}
\emph{Existence}: by composition of
Theorems~\ref{thm:physical-hilbert-space}, \ref{thm:semigroup}, and
\ref{thm:OS-hamiltonian}.  \emph{Uniqueness up to unitary
equivalence}: the construction is canonical given $(\muinf, \Ecalp,
\theta, \Tt)$; any two reconstructions
$\mathrm{Rec}(\muinf)$ and $\mathrm{Rec}'(\muinf)$ have isometric
$\Ecalp / \Ncal$ pre-Hilbert spaces (by the same Schwinger inner
product), hence isometric completions, hence unitarily equivalent
Hamiltonians.
\end{proof}

\section{Preparation for the mass-gap chapter: the spectral injection}
\label{sec:phase-9-prep}

\begin{proposition}[Spectral content of $\Hop$ and the mass gap]
\label{prop:spectral-mass-gap}
The spectrum $\sigma(\Hop) \subseteq [0, \infty)$ decomposes as
\begin{equation*}
  \sigma(\Hop) \;=\; \{0\} \;\cup\; \sigma_+(\Hop),
\end{equation*}
where $\sigma_+(\Hop) := \sigma(\Hop) \cap (0, \infty)$ is the
\emph{excited spectrum}.  The vacuum $\Omegavac$ spans the
$0$-eigenspace.  The \emph{mass gap} is
\begin{equation*}
  m \;:=\; \inf \, \sigma_+(\Hop) \;\in\; [0, \infty].
\end{equation*}
The mass-gap statement is $m > 0$, equivalently
$\sigma(\Hop) \cap (0, m) = \emptyset$ for some $m > 0$
(Clay-statement form, the work plan \S9).
\end{proposition}

\begin{proof}
Standard spectral theorem decomposition for self-adjoint operators;
see Reed--Simon Vol.~I Chapter~VII.  Vacuum-annihilation
($\Hop \Omegavac = 0$) gives the $0$-eigenspace; the rest of the
spectrum is the excited spectrum.
\end{proof}

\begin{theorem}[OS4 cluster property $\Rightarrow$ spectral gap]
\label{thm:cluster-to-gap}
If $\muinf$ satisfies OS4 with explicit exponential cluster rate
$m_0 > 0$ (\texttt{os-axioms.tex} Theorem~7.1; the rate is
$m_0 := \sighinf$ from the $\sighinf$-positivity analysis), then the OS Hamiltonian satisfies
\begin{equation*}
  \sigma(\Hop) \cap (0, m_0) \;=\; \emptyset,
  \qquad \text{i.e., } \;\; m \;\ge\; m_0 \;>\; 0.
\end{equation*}
\end{theorem}

\begin{proof}[Proof sketch]
The two-point Schwinger function under $\muinf$ has the spectral
representation
\begin{equation*}
  \int (\theta F)^* (\Tt G) \, d\muinf
  \;=\;
  \IprodOS{F}{\Ut G}
  \;=\;
  \int_0^\infty e^{-tE} \, d\mu_{F,G}(E),
\end{equation*}
where $d\mu_{F,G}(E)$ is the spectral measure of $\Hop$ between
$[F]$ and $[G]$ in $\Hphys$.  OS4 says the LHS decays as
$\le C \, e^{-m_0 \cdot t}$ for spatially-separated $F, G$ at large
$t$; the RHS decay rate is the lower edge of $\mathrm{supp} \, d\mu_{F,G}$.
Combining: $\inf \mathrm{supp} \, d\mu_{F,G} \ge m_0$ for every
$F, G$ with $\IprodOS{F}{\Omegavac} = \IprodOS{G}{\Omegavac} = 0$
(i.e., orthogonal to the vacuum).  Taking the supremum over such
$F, G$ gives $\inf \sigma_+(\Hop) \ge m_0$.

The detailed argument is Glimm--Jaffe \S6.1 (the
``mass-gap-from-clustering'' theorem); see also Simon
\emph{The $P(\phi)_2$ Euclidean Quantum Field Theory} \S6.
\end{proof}

Theorem~\ref{thm:cluster-to-gap} is the analytical bridge from OS4
of the OS-axiom verification to the spectral gap of the mass-gap chapter.  Its precondition
(OS4 with explicit rate $m_0 > 0$) is the $\sighinf$-positivity analysis + 8.1; its
conclusion ($m \ge m_0 > 0$) is the SU(2) mass-gap statement at the
Hilbert-space level.

\section{Lean type signature and scope}
\label{sec:lean-signature}

The Lean module \texttt{lean/YMGap/Continuum/OSReconstruction.lean}
exposes:
\begin{itemize}
  \item \texttt{HilbertSpace} — abstract type for the physical
    Hilbert space.
  \item \texttt{SelfAdjointOp} — abstract type for the Hamiltonian
    operator on \texttt{HilbertSpace}.
  \item \texttt{VacuumVector} — abstract type for the
    distinguished vacuum.
  \item \texttt{OSReconstruction} — the structure bundling
    $(\Hphys, \Omegavac, \Hop)$ together with the abstract
    semigroup data.
  \item \texttt{os\_reconstruct\_su2} — smart constructor taking
    the paper-level Prop inputs (existence of $\Hphys$, etc.) and
    assembling the reconstruction.
  \item \texttt{cluster\_implies\_gap} — Prop signature for
    Theorem~\ref{thm:cluster-to-gap}, the the mass-gap chapter mass-gap input.
\end{itemize}

\subsection*{Decision: paper-only existence + Lean-formal-statement-only}

Per the work-plan \Cref{ch:os-reconstruction} open question (``Lean formalization
of the reconstruction is a major undertaking [...] likely deferred
to paper-first, Lean-later or a parallel Lean-formal-statement-only
milestone''), this project commits to the \textbf{paper-first +
Lean-formal-statement-only} branch.  Rationale:

\begin{enumerate}
  \item Stone's theorem + the spectral theorem for unbounded
    self-adjoint operators require a major constructive-analysis
    Lean build-out (densely-defined operators, spectral measures,
    functional calculus) that is far beyond the current Q'-rational
    + projective-system stack.
  \item The the mass-gap chapter mass-gap statement can be expressed against the
    \texttt{OSReconstruction} type signature; the paper-level
    existence proof here + the Theorem~\ref{thm:cluster-to-gap}
    bridge lets the mass-gap chapter land in Lean without waiting on Stone's
    theorem.
  \item a constructive-reals layer+ reserves the option of building the
    constructive analysis layer later (with $\Hop$ as one motivating
    use case); this decision does not foreclose that.
\end{enumerate}

This is parallel to the the continuum-limit analysis.2 decision recorded in
the supporting analysis and the the OS-axiom verification decision
recorded in the supporting analysis  Each of
the three milestones (7.2, 8.1, 8.2) explicitly opts for the
paper-only / Lean-formal-statement-only branch with the same
rationale.

\section{Open analytical items}
\label{sec:open}

\begin{itemize}
  \item \textbf{Lean formalization of Stone's theorem.}  Required
    to refine the abstract \texttt{SelfAdjointOp} type into a Lean
    object with a real domain, spectral measure, functional
    calculus.  Major undertaking; deferred.
  \item \textbf{Sharper $m \ge m_0$ bound.}
    Theorem~\ref{thm:cluster-to-gap} gives $m \ge m_0$ where
    $m_0 = \sighinf$ from the $\sighinf$-positivity analysis; the inequality may be strict
    (the actual mass could be larger than the σ̂-derived lower bound).
    Sharpness is a the mass-gap chapter / the SU(N) extension framework open analytic item.
  \item \textbf{Stability under perturbations.}  The OS reconstruction
    is stable under small perturbations of $\muinf$ (Hille--Yosida
    perturbation theory); making this constructive at the rate
    needed for the Phase-10 SU(3) extension is downstream work.
  \item \textbf{Operator-valued distribution structure.}  The full
    Wightman reconstruction of Streater--Wightman fields from the
    Schwinger functions is a further analytical step, not addressed
    here; we deliver the Hilbert space + Hamiltonian, which is what
    the the mass-gap chapter mass-gap statement needs.
\end{itemize}

\section*{References}

\noindent[GJ87] J.~Glimm, A.~Jaffe, \emph{Quantum Physics: A
Functional Integral Point of View}, 2nd ed., Springer 1987.
\S6.1: OS axioms and the reconstruction theorem (the canonical
reference for the construction here).

\noindent[OS73] K.~Osterwalder, R.~Schrader, ``Axioms for Euclidean
Green's functions,'' Comm.\ Math.\ Phys.\ \textbf{31} (1973),
83--112.  Original OS reconstruction paper.

\noindent[OS75] K.~Osterwalder, R.~Schrader, ``Axioms for Euclidean
Green's functions.~II,'' Comm.\ Math.\ Phys.\ \textbf{42} (1975),
281--305.  Refinement and clustering-to-mass-gap argument.

\noindent[RS-II] M.~Reed, B.~Simon, \emph{Methods of Modern
Mathematical Physics}, Vol.~II: Fourier Analysis, Self-Adjointness,
Academic Press 1975.  Theorems X.47a and X.49 on contraction
semigroups + Stone's theorem.

\noindent[Sim74] B.~Simon, \emph{The $P(\phi)_2$ Euclidean (Quantum)
Field Theory}, Princeton 1974.  \S6: clustering-to-mass-gap.

\noindent[Jor86] P.~Jorgensen, ``Analytic continuation of local
representations of Lie groups,'' Pac.\ J.\ Math.\ \textbf{125}
(1986), 397--408.  Operator-valued analytic continuation,
OS-positivity from a Lie-group representation-theoretic angle.

\noindent[JW06] A.~Jaffe, E.~Witten, ``Quantum Yang--Mills theory,''
Clay Millennium Problem statement (2006).
